\chapter{Introduction}

\section{Purpose}
BBP (Bike Path Platform) is designed to address the challenges cyclists face when searching for safe and well-maintained bike routes. The platform enables users to record their cycling trips, either automatically through device sensors or manually by inputting street names, and to share this information with the community. By aggregating data from multiple users, BBP builds a comprehensive database of bike paths enriched with real-time conditions, obstacle information, and community-validated status.

The system aims to provide cyclists ranked path recommendations based on aggregated community feedback, detected road anomalies, and current weather conditions. Additionally, the platform supports obstacle detection through sensor data analysis, allowing users to validate and report road hazards that are then integrated into the public path network.

\subsection{Goals}
\begin{itemize}
    \item \textbf{[G0]:} A user wants to signup first and then login to the application.
    \item \textbf{[G1]:} A registered user aims to automatically record their cycling trips using theirs mobile device's sensors. 
    \item \textbf{[G2]:} A registered user aims to manually record a bike trip by inserting one street at a time without physical cycling.
    \item \textbf{[G3]:} A registered user wants to access their recorded bike trips history.
    \item \textbf{[G4]:} A registered user aims to view trip on map and consult its statistics, including average speed, trip time and meteorological conditions (at recording time).
    \item \textbf{[G5]:} A registered user wants to share information about a bike trip they have recorded, including their status and the presence of obstacles, making this information available to the community (all users).
    \item \textbf{[G6]:} A User wants to search for bike paths given an origin and a destination point and visualize results on a map.
    \item \textbf{[G7]:} A User wants to retrieve up-to-date information about bike paths, their status, and current weather conditions.
    \item \textbf{[G8]:} A registered user wants to delete one of their previously recorded trips.
\end{itemize}

\section{Scope}
The BBP platform provides a comprehensive set of features to support the cycling community in finding safe routes and contributing to collective knowledge about urban bike paths.

Users can register and authenticate through the mobile application to access all platform features. The core functionality revolves around trip recording, which can be performed in two modes: automatic recording leverages device sensors (GPS, accelerometer, gyroscope) to track the user's movement and detect potential road anomalies such as potholes, while manual recording allows users to input street names to define their route without physical cycling.

After completing an automatic trip, the system performs asynchronous analysis to compute statistics (distance, duration, average speed) and identify obstacle candidates. Users must validate these detected obstacles before publishing their trip. Once validated, trips can be published to contribute anonymized data to the public path network.

The platform aggregates contributions from multiple users through a merging process that computes path scores and updates road status. Users can search for paths between any origin and destination, with results ranked by score and enriched with current weather information retrieved from external services.

The architecture follows a microservices-based design with event-driven communication, ensuring scalability, loose coupling, and resilience. External Map and Weather services are integrated through adapters protected by circuit breaker mechanisms to handle third-party unavailability gracefully.

\section{Definitions, Acronyms, Abbreviations}

\subsection{Definitions}
\begin{itemize}
    \item \textbf{Registered user:} User that is registered in the platform.
    \item \textbf{Guest user:} User that is not registered in the platform.
    \item \textbf{User:} Term used to indicate both registered and guest users.
    \item \textbf{Path:} Public route connecting two geographical points, obtained from the merging process of multiple Trips sharing identical GPS coordinates and cumulative distance metrics. This designation indicates either the presence of dedicated cycling infrastructure or low-traffic roadways with speed limits compatible with typical cycling velocities.
    \item \textbf{Trip:} A user-specific ordered sequence of GPS coordinates or street identifiers, accompanied by associated meteorological data, recorded by a registered user.
    \item \textbf{Merging process:} A computational process through which the BBP system aggregates trip data, potentially coming from multiple users, to generate public paths. It also represents the process of updating an already-existing public path with recently published trips.
    \item \textbf{Anomaly:} Pothole or obstacle found while recording the trip and identified by the sensors data analysis.
    \item \textbf{Trip status:} The maintenance condition detected or validated for a user trip, based on sensor data analysis (potholes) or manual insertion.
    \item \textbf{Path status:} The aggregated maintenance condition of a bike path, computed by the BBP system whenever a new trip is published and merged. The calculation considers the status of all merged trips, weighing factors such as data freshness and consensus among multiple users' reports (e.g., if two trips report ``optimal'' and one reports ``requires maintenance'' within the same timeframe, the path is considered ``optimal'').
    \item \textbf{Path score:} A quality metric used to rank and visualize bike paths, calculated based on the path's status and its effectiveness in connecting origin to destination (measured through average trip duration).
    \item \textbf{Corresponding Path (to a Trip):} The path that shares the same trajectory (ordered set of GPS coordinates) as a given trip.
    \item \textbf{Trip Merging Timeframe:} The time window within which trip reports are considered for consensus-based status aggregation during the path merging process.
\end{itemize}

\subsection{Acronyms}
\begin{itemize}
    \item \textbf{API:} Application Programming Interface
    \item \textbf{BBP:} Best Bike Paths
    \item \textbf{GPS:} Global Positioning System
    \item \textbf{JWT:} JSON Web Token
    \item \textbf{DBMS:} DataBase Management System
    \item \textbf{TLS:} Transport Layer Security
    \item \textbf{HTTPS:} Hypertext Transfer Protocol Secure
    \item \textbf{CRUD:} Create, Read, Update, Delete
    \item \textbf{DTO:} Data Transfer Object
\end{itemize}

\subsection{Abbreviations}
\begin{itemize}
    \item \textbf{DB:} Database
    \item \textbf{LB:} Load Balancer
    \item \textbf{Auth:} Authentication and Authorization
    \item \textbf{Pub/Sub:} Publish/Subscribe
\end{itemize}

\section{Revision History}
\begin{table}[H]
\centering
\begin{tabular}{|l|l|p{8cm}|}
\hline
\textbf{Version} & \textbf{Date} & \textbf{Description} \\ \hline
1.0 & 06/01/2025 & First Release \\ \hline
\end{tabular}
\caption{Revision History}
\label{tab:revision_history}
\end{table}

\section{Reference Documents}
\begin{itemize}
    \item Project assignment document
    \item Software Engineering 2 Course Material - A.Y. 2025/2026
    \item Requirements Analysis and Specification Document for the BBP platform
\end{itemize}

\section{Document Structure}
The document is organized into the following chapters:

\begin{table}[H]
    \centering
    % Imposta la tabella larga quanto il testo (\textwidth)
    % La colonna X si adatta automaticamente allo spazio rimasto
    \begin{tabularx}{\textwidth}{|l|p{4cm}|X|}
    \hline
    \textbf{Chapter} & \textbf{Title} & \textbf{Description} \\ \hline
    C1 & Introduction & Introduces the goals and purpose of the project, defines terminology, acronyms, and abbreviations needed to fully comprehend the document \\ \hline
    C2 & Architectural Design & Describes the BBP platform architecture at multiple levels of granularity, including component view, deployment view, runtime view (sequence diagrams), component interfaces, and architectural styles and patterns \\ \hline
    C3 & User Interface Design & Provides visual representations of the mobile application user interface for all main user flows \\ \hline
    C4 & Requirements Traceability & Maps the functional requirements defined in the RASD to the specific architectural components and subsystems designed in this document \\ \hline
    C5 & Implementation, Integration and Testing Plan & Outlines the bottom-up implementation strategy, the order in which components are developed, tested, and integrated, and the testing approach at each level \\ \hline
    C6 & Effort Spent & Documents the workload distribution among team members \\ \hline
    C7 & References & Lists the external resources and tools used to create this document \\ \hline
    \end{tabularx}
    \caption{Document Structure}
    \label{tab:document_structure}
\end{table}